\section{Pressure-dependent dunite experiment}

Lawal and Kim \citep{lawal2026} performed drained and unjacketed hydrostatic
compression tests on a dry heat-treated control (HT14) and specimens
serpentinized for 14 and 30 days (SP14 and SP30).  They report
\begin{equation}
 \alpha(P)=1-\frac{K(P)}{K_s'},
 \label{eq:lawal-kim-alpha}
\end{equation}
where \(P\) is confining pressure, \(K(P)\) is the drained tangent bulk modulus,
and \(K_s'\) is the unjacketed modulus.  The public workbook contains 46 plotted
pressure--modulus--coefficient points.  Recalculation of
\cref{eq:lawal-kim-alpha} reproduces every stored coefficient to machine
precision.

We represent compliant crack closure with
\begin{equation}
 \frac{1}{K(P)}=\frac{1}{K_\infty}+C_c\exp\left(-\frac{P}{P_c}\right),
 \label{eq:crack-compliance}
\end{equation}
where \(K_\infty>0\) is the high-pressure drained modulus, \(C_c>0\) is the
initial crack compliance, and \(P_c>0\) is a closure pressure.  Positivity of
these parameters guarantees a positive tangent modulus.

For the nonreacting constitutive path, pressure is measured in GPa and the local
constraints are
\begin{align}
 0={}&\ln J_s+\frac{P}{K_\infty}
 +C_cP_c\left[1-\exp\left(-\frac{P}{P_c}\right)\right],
 \label{eq:hydrostatic-constraint}\\
 0={}&J_s\phi_s-\exp\left(-\frac{P}{K_s'}\right).
 \label{eq:grain-volume-constraint}
\end{align}
The first equation integrates the fitted drained tangent compliance.  The
second supplies a compressible intrinsic grain volume under unjacketed loading,
with normalized reference specific volume and fixed referential solid
accumulation.  Implicit differentiation gives
\begin{equation}
 \frac{\partial P}{\partial J_s}=-\frac{K(P)}{J_s},
 \qquad
 \left.\frac{\partial\bar v_s}{\partial J_s}\right|_{p_E}
 =\frac{\bar v_s}{J_s}\frac{K(P)}{K_s'},
 \label{eq:pressure-path-tangent}
\end{equation}
and therefore
\begin{equation}
 B(P)=1-\frac{\bar v_s}{J_s v_{s0}}\frac{K(P)}{K_s'}.
 \label{eq:finite-pressure-biot}
\end{equation}
Equation~\eqref{eq:finite-pressure-biot} reduces to the experimental expression
\eqref{eq:lawal-kim-alpha} as \(\bar v_s/(J_sv_{s0})\to1\).  It also provides a
small finite-deformation correction rather than imposing the measured
coefficient directly.

The fit minimizes unweighted residuals in the measured drained modulus.  No
coefficient values are used independently because the released coefficients
are algebraically generated from the same modulus data.  The workbook uses
\(K_s'=65.0\)~GPa in the SP14 formulas, whereas the annotation in the published
Figure~3b reads \(65.5\)~GPa.  We use 65.0~GPa to reproduce Figure~3c exactly
and retain this discrepancy in the data-provenance audit.

